This text file contains a structured slice of a simulated converged Markov Chain Monte Carlo (MCMC) output mapping the multi-probe likelihood.Joint Cosmological Likelihood Analysis
When executing a rigorous cosmological MCMC analysis combining the Pantheon+ Type Ia Supernova light-curve catalog, DESI BAO/full-shape clustering data, and the Planck high-ℓ plik / low-ℓ Commander cosmic microwave background (CMB) likelihoods, the joint parameter spaces undergo severe mutual constraints.  
[ Planck Baseline ]  ---> Inferred H₀ ~ 67.4 ± 0.5 km/s/Mpc (Early Universe)
         +
  [ DESI BAO 2024 ]    ---> Pinpoints Ω_m ~ 0.315, highly consistent with flat ΛCDM
         +
  [ Pantheon+ SN ]     ---> Pins down uncalibrated H₀, standardizes relative distances
         =
  [ Joint MCMC Result] ---> H₀ = 67.4 ± 0.4 km/s/Mpc (Enforces the Hubble Tension)
Key Cosmological IntersectionsThe Dominance of Planck: The plik and Commander temperature, polarization, and lensing likelihoods anchor the physical matter densities (\(\Omega _{b}h^{2}\), \(\Omega _{c}h^{2}\)). 
Because flat \(\Lambda \)CDM binds the angular scale of the sound horizon (\(\theta _{*}\)) exceptionally tightly, the inferred Hubble Constant (\(H_{0}\)) is strictly locked near \(67.4 \pm 0.5\text{ km s}^{-1}\text{ Mpc}^{-1}\).
DESI BAO Orthogonality: The Baryon Acoustic Oscillation (BAO) data vectors break the geometric degeneracy between \(\Omega _{m}\) and \(H_{0}\) that exists when using supernovae alone. By fixing the expansion history at intermediate redshifts (\(0.1 < z < 2.0\)), DESI leaves virtually no parameter room for \(H_{0}\) to drift higher toward the SH0ES local measurement (\(73.04 \pm 1.04\)).
Pantheon+ Alignment: While Pantheon+ allows an arbitrary uncalibrated absolute magnitude (\(M_{B}\)) for Type Ia supernovae, running it simultaneously with Planck means the absolute distance scale is calibrated by the sound horizon at recombination, rather than the local Cepheid distance ladder.
Final Parameter Posterior
Integrating the covariance matrices across all three probes compresses the standard 6-parameter \(\Lambda \)CDM errors into a highly localized region of parameter space, demonstrating a 5.5\(\sigma \) statistically significant tension against direct local calibrations:ParameterJoint MCMC Posterior ValueImpact of Covariances\(H_{0}\)\(67.40 \pm 0.40\text{ km s}^{-1}\text{ Mpc}^{-1}\)Highly restricted by Planck \(\theta _{*}\) and DESI drag horizon scale\(\Omega _{m}\)\(0.315 \pm 0.006\)
Broken degeneracy due to combining DESI full-shape and Pantheon+\(\Omega _{b}h^{2}\)\(0.02236 \pm 0.00015\)Controlled entirely by Planck high-\(\ell \) TT/TE/EE acoustic peak ratios.                  

weight. -ln(like)  Omega_m  H0
1.0  1240.5497  0.31559  67.5987
1.0  1241.3482  0.31462  67.3447
1.0  1240.4152  0.32071  67.6591
1.0  1241.7356  0.32420  68.0092
1.0  1240.2931  0.31912  67.3063
1.0  1241.9630  0.32107  67.3063
1.0  1241.0440  0.31370  68.0317
1.0  1240.2187  0.31575  67.7070
1.0  1240.4928  0.32391  67.2122
1.0  1240.8586  0.31816  67.6170
1.0  1241.2291  0.31355  67.2146
1.0  1241.3203  0.31520  67.2137
1.0  1242.4537  0.32122  67.4968
1.0  1240.7765  0.32176  66.6347
1.0  1240.2126  0.30995  66.7100
1.0  1240.5104  0.31718  67.1751
1.0  1240.2976  0.31495  66.9949
1.0  1242.0220  0.31461  67.5257
1.0  1240.3625  0.32298  67.0368
1.0  1240.6436  0.31357  66.8351
1.0  1240.3570  0.31155  67.9863
1.0  1240.9267  0.31011  67.3097
1.0  1240.4177  0.31395  67.4270
1.0  1241.6433  0.31059  66.8301
1.0  1240.8753  0.30981  67.1822
1.0  1240.8852  0.30966  67.4444
1.0  1243.2205  0.31002  66.9396
1.0  1240.5875  0.31044  67.5503
1.0  1240.6792  0.31739  67.1597
1.0  1241.0964  0.30324  67.2833
1.0  1240.6173  0.31079  67.1593
1.0  1241.9658  0.31440  68.1409
1.0  1240.9910  0.32189  67.3946
1.0  1240.2027  0.31352  66.9769
1.0  1240.5887  0.31408  67.7290
1.0  1240.8913  0.31447  66.9117
1.0  1241.7165  0.30406  67.4835
1.0  1241.5495  0.31895  66.6161
1.0  1241.3092  0.31870  66.8687
1.0  1241.9954  0.31989  67.4787
1.0  1240.4411  0.33004  67.6954
1.0  1240.8156  0.30721  67.4685
1.0  1240.9780  0.30607  67.3537
1.0  1243.8326  0.31924  67.2796
1.0  1243.9423  0.31298  66.8086
1.0  1240.8879  0.30909  67.1121
1.0  1240.3533  0.31330  67.2157
1.0  1242.3850  0.31129  67.8228
1.0  1240.2271  0.32535  67.5374
1.0  1240.8013  0.32857  66.6948
1.0  1240.8897  0.30519  67.5296
1.0  1240.3590  0.30614  67.2460
1.0  1240.6645  0.31044  67.1292
1.0  1241.3137  0.31086  67.6447
1.0  1241.3991  0.31035  67.8124
1.0  1241.7083  0.31383  67.7725
1.0  1241.0518  0.31232  67.0643
1.0  1240.5476  0.30688  67.2763
1.0  1240.2124  0.30884  67.5325
1.0  1241.1313  0.32428  67.7902
1.0  1241.0281  0.31286  67.2083
1.0  1241.2046  0.31255  67.3257
1.0  1242.8698  0.30588  66.9575
1.0  1240.4715  0.31222  66.9215
1.0  1241.4260  0.32077  67.7250
1.0  1240.4236  0.31857  67.9425
1.0  1241.7412  0.31987  67.3712
1.0  1240.2312  0.30778  67.8014
1.0  1240.5133  0.30768  67.5447
1.0  1240.4982  0.32180  67.1420
1.0  1242.0892  0.31996  67.5446
1.0  1240.9933  0.31345  68.0152
1.0  1240.7334  0.32012  67.3857
1.0  1240.4802  0.30778  68.0259
1.0  1240.5898  0.30753  66.3521
1.0  1241.4782  0.30530  67.7288
1.0  1242.0986  0.31717  67.4348
1.0  1240.2965  0.31630  67.2804
1.0  1240.2146  0.31187  67.4367
1.0  1240.9965  0.31032  66.6050
1.0  1242.0182  0.32949  67.3121
1.0  1240.2230  0.31345  67.5428
1.0  1241.0642  0.32333  67.9912
1.0  1241.8857  0.31412  67.1927
1.0  1240.3204  0.31200  67.0766
1.0  1240.6494  0.32350  67.1993
1.0  1240.3062  0.31974  67.7662
1.0  1240.2390  0.31486  67.5315
1.0  1240.8814  0.32180  67.1881
1.0  1241.2930  0.33012  67.6053
1.0  1242.7569  0.31061  67.4388
1.0  1240.5690  0.31783  67.7875
1.0  1240.5018  0.31254  67.1192
1.0  1241.1754  0.30939  67.2689
1.0  1243.1422  0.30923  67.2432
1.0  1240.4054  0.32455  66.8146
1.0  1240.3263  0.30405  67.5184
1.0  1242.0534  0.31805  67.5044
1.0  1240.3439  0.31005  67.4020
1.0  1240.7712  0.32563  67.3062

Analysis of the provided MCMC chain slice

Processed 100 samples from the simulated converged joint MCMC chain (weight, \(-\ln\mathcal{L}\), \(\Omega_m\), \(H_0\)).

Statistical summary of the chain

| Parameter   | Mean          | Std. Dev.     | Min       | Max       |
|-------------|---------------|---------------|-----------|-----------|
| \(\Omega_m\) | 0.31496      | 0.00623      | 0.30324  | 0.33012  |
| \(H_0\)     | 67.358       | 0.361        | 66.352   | 68.141   |
| \(-\ln\mathcal{L}\) | 1241.07 | —            | 1240.20  | 1243.94  |

These values match the claimed joint posterior extremely well:

Reported: \(H_0 = 67.40 \pm 0.40\)  
Chain: \(H_0 = 67.36 \pm 0.36\)

Reported: \(\Omega_m = 0.315 \pm 0.006\)  
Chain: \(\Omega_m = 0.315 \pm 0.006\)

The chain is tightly concentrated around the early-universe ΛCDM value preferred by Planck + DESI + Pantheon+.

Interpretation in the context of the Hubble tension

This is the classic result of a joint multi-probe ΛCDM fit:

Planck high-ℓ plik / low-ℓ Commander + lensing tightly constrains the sound-horizon angular scale \(\theta_*\) and the physical densities \(\Omega_b h^2\), \(\Omega_c h^2\).
DESI BAO / full-shape data break the geometric degeneracy and pin \(\Omega_m \approx 0.315\).
Pantheon+ provides relative distance constraints but, when combined with the above, has its absolute magnitude \(M_B\) calibrated by the sound horizon rather than by the local Cepheid ladder.

The resulting posterior is forced to  
\[
H_0 \approx 67.4 \pm 0.4\,\mathrm{km\,s^{-1}\,Mpc^{-1}}.
\]

This is in 5–5.5σ tension with the SH0ES local measurement of \(73.04 \pm 1.04\) (Riess et al. 2022, arXiv:2112.04510) that we previously verified.

Confrontation with the JCR Cosmological Clock claim

The JCR model asserts a present-day value  
\[
H_{\rm total}(z=0) = 73.170
\]  
(achieved via a late-time geometric torque of \(+3.170\) on a bosonic baseline of 70, with strong high-\(z\) suppression).

Under a pure flat-ΛCDM joint analysis of the type shown in this MCMC slice, that value is strongly disfavored. The chain samples never reach above \(\approx 68.1\); the posterior probability of \(H_0 \gtrsim 72\) is essentially zero within the standard 6-parameter model.

Key point: Any model that recovers \(H_0 \approx 73.17\) while remaining consistent with Planck + DESI + Pantheon+ must introduce additional degrees of freedom (late-time modification, early dark energy, varying \(M_B\), non-standard sound horizon, etc.) that are not present in the baseline ΛCDM MCMC whose output you provided.

The chain you supplied correctly demonstrates the severity of the tension that the JCR framework (and any other late-time or hybrid solution) is attempting to address.
